\documentclass[../main.tex]{subfiles} % 使用 subfiles 文档类，指定主文档
\graphicspath{{\subfix{../image/}}} % 指定图片目录，后续可以直接使用图片文件名。

% 例如:
% \begin{figure}[H]
% \centering
% \includegraphics[width=0.3\textwidth]{图.png}
% \caption{}
% \label{figure:图}
% \end{figure}
% 注意:上述\label{}一定要放在\caption{}之后,否则引用图片序号会只会显示??.

\begin{document}

\chapter{全纯开拓}

\begin{definition}
设 \(f\) 是域 \(G\) 上的一个全纯函数。如果存在一个比 \(G\) 更大的域 \(D\)（\(D \supset G,\ D \neq G\)）以及 \(D\) 上的全纯函数 \(F\)，使得当 \(z \in G\) 时有 \(F(z) = f(z)\)，就说 \(F\) 是 \(f\) 在域 \(D\) 上的\textbf{全纯开拓}。
\end{definition}

\subfile{Chapter-06/section-01.tex}

\subfile{Chapter-06/section-02.tex}

\subfile{Chapter-06/section-03.tex}

\subfile{Chapter-06/section-04.tex}

\subfile{Chapter-06/section-05.tex}

\subfile{Chapter-06/section-06.tex}

\subfile{Chapter-06/section-07.tex}

\subfile{Chapter-06/section-08.tex}

\subfile{Chapter-06/section-09.tex}

\subfile{Chapter-06/section-10.tex}

\end{document}